%----------Anhang/Appendices--------------------------------------------------------------

\appendix
% =========================================================
% APPENDICES
% =========================================================

\section{Appendix A: Repository Layout and Artifact Naming}
\label{app:repo_layout}

A consistent repository layout is used to separate source code, configuration, prompts/schemas, datasets, and run outputs. This structure supports prompt/schema versioning, deterministic replay of stages, stage-wise evaluation (Section~7), and integration into batch workflows (Section~4).

\subsection{Recommended repository structure}
\begin{verbatim}
DocValidator/
  optimal/
    baseline_check.py
    cosine_tfidf_layer.py
    llm_preprocess.py
    llm_pass1.py
    llm_pass2.py
    pipeline_runner.py

  schemas/
    pass1_schema.json
    pass2_schema.json
    preprocessed_schema.json
    baseline_report_schema.json
    cosine_report_schema.json

  prompts/
    pass1_system.txt
    pass2_system.txt

  config/
    thresholds.yaml
    pipeline.yaml
    logging.yaml

  data/
    template/
      expected_spec.json
    parsed/
      <doc_id>.json
    reports/
      baseline/      baseline_report_<doc_id>.json
      cosine/        cosine_resolved_report_<doc_id>.json
      preprocessed/  preprocessed_<doc_id>.json
      pass1/         pass1_audit_<doc_id>.json
      pass2/         pass2_consistency_<doc_id>.json

  runs/
    <run_id>/
      run_manifest.json
      logs/
      outputs/

  README.md
\end{verbatim}

\subsection{Artifact naming conventions}
Intermediate artifacts are persisted intentionally to enable debugging and stage contribution analysis:
\begin{itemize}
    \item \texttt{baseline\_report\_<doc\_id>.json}
    \item \texttt{cosine\_resolved\_report\_<doc\_id>.json}
    \item \texttt{preprocessed\_<doc\_id>.json}
    \item \texttt{pass1\_audit\_<doc\_id>.json}
    \item \texttt{pass2\_consistency\_<doc\_id>.json}
\end{itemize}


\section{Appendix B: Configuration Files (Thresholds and Pipeline Execution)}
\label{app:config_files}

Thresholds and stage toggles are externalized to configuration files. Externalizing parameters avoids hard-coded ``magic constants'' and enables controlled tuning and sensitivity analysis in the evaluation.

\subsection{Threshold configuration (example)}
\begin{verbatim}
# config/thresholds.yaml
baseline:
  empty_min_chars: 40
  prose_min_chars_for_dup_check: 150

cosine:
  section_rescue_threshold: 0.55
  placeholder_heading_threshold: 0.95
  guidance_copy_threshold: 0.99
  duplicate_threshold: 0.98
  legal_note_min_similarity: 0.85
\end{verbatim}

\subsection{Pipeline execution configuration (example)}
\begin{verbatim}
# config/pipeline.yaml
stages:
  run_baseline: true
  run_cosine: true
  run_preprocess: true
  run_pass1: true
  run_pass2: true

paths:
  template_spec: "data/template/expected_spec.json"
  parsed_dir: "data/parsed"
  baseline_dir: "data/reports/baseline"
  cosine_dir: "data/reports/cosine"
  preprocessed_dir: "data/reports/preprocessed"
  pass1_dir: "data/reports/pass1"
  pass2_dir: "data/reports/pass2"
  runs_dir: "runs"
\end{verbatim}


\section{Appendix C: Prompt and Schema Versioning}
\label{app:prompt_schema_versioning}

Prompts and JSON schemas are treated as first-class versioned artifacts to ensure that audit outputs remain comparable across runs.

\subsection{Prompt files}
\begin{itemize}
    \item \texttt{prompts/pass1\_system.txt}: integrity and semantic adequacy auditing (evidence-only, no unsupported claims).
    \item \texttt{prompts/pass2\_system.txt}: entity extraction and cross-view consistency checks (BBV$\leftrightarrow$Runtime, BBV$\leftrightarrow$Deployment, Goals$\leftrightarrow$Decisions, Decisions$\leftrightarrow$Risks).
\end{itemize}

\subsection{Schema files}
\begin{itemize}
    \item \texttt{schemas/pass1\_schema.json}: structured output for Pass~1.
    \item \texttt{schemas/pass2\_schema.json}: structured output for Pass~2.
    \item \texttt{schemas/preprocessed\_schema.json}: validation schema for the preprocessed input representation (recommended).
\end{itemize}


\section{Appendix D: Run Metadata, Run Manifest, and CLI Specification}
\label{app:run_metadata_cli}

\subsection{Run metadata embedded in outputs}
Each output artifact includes a \texttt{run\_metadata} block to support traceability and reproducibility:

\begin{verbatim}
{
  "run_metadata": {
    "run_id": "2026-01-14_01",
    "tool_version": "git:<commit-hash>",
    "template_id": "company_arc42_variant_vX",
    "template_hash": "<sha256>",
    "thresholds_file": "config/thresholds.yaml",
    "thresholds_hash": "<sha256>",
    "prompt_pass1_hash": "<sha256>",
    "prompt_pass2_hash": "<sha256>",
    "schema_pass1_hash": "<sha256>",
    "schema_pass2_hash": "<sha256>",
    "model_pass1": "<model-name>",
    "model_pass2": "<model-name>",
    "timestamp_utc": "2026-01-14T18:21:55Z"
  }
}
\end{verbatim}

\subsection{Run manifest per execution}
A run manifest summarizes inputs, configuration, and outputs for a given run:

\begin{verbatim}
{
  "run_id": "2026-01-14_01",
  "timestamp_utc": "2026-01-14T18:21:55Z",
  "tool_version": "git:<commit-hash>",
  "inputs": {
    "template_spec": "data/template/expected_spec.json",
    "parsed_dir": "data/parsed",
    "parsed_snapshot_hash": "<sha256>"
  },
  "configuration": {
    "pipeline_yaml": "config/pipeline.yaml",
    "thresholds_yaml": "config/thresholds.yaml",
    "thresholds_hash": "<sha256>"
  },
  "prompts_and_schemas": {
    "prompt_pass1": "prompts/pass1_system.txt",
    "prompt_pass2": "prompts/pass2_system.txt",
    "prompt_pass1_hash": "<sha256>",
    "prompt_pass2_hash": "<sha256>",
    "schema_pass1": "schemas/pass1_schema.json",
    "schema_pass2": "schemas/pass2_schema.json",
    "schema_pass1_hash": "<sha256>",
    "schema_pass2_hash": "<sha256>"
  },
  "stages_executed": {
    "baseline": true,
    "cosine": true,
    "preprocess": true,
    "pass1": true,
    "pass2": true
  },
  "counts": {
    "documents_processed": 0,
    "documents_failed": 0
  },
  "outputs": {
    "baseline_dir": "data/reports/baseline",
    "cosine_dir": "data/reports/cosine",
    "preprocessed_dir": "data/reports/preprocessed",
    "pass1_dir": "data/reports/pass1",
    "pass2_dir": "data/reports/pass2"
  }
}
\end{verbatim}

\subsection{Command-line interface (CLI) specification}
A unified runner script executes the pipeline deterministically.

\paragraph{Usage}
\begin{verbatim}
python optimal/pipeline_runner.py --config config/pipeline.yaml --run-id 2026-01-14_01
\end{verbatim}

\paragraph{Arguments}
\begin{itemize}
    \item \texttt{--config <path>} (required): path to \texttt{pipeline.yaml}.
    \item \texttt{--run-id <string>} (optional): unique run identifier; defaults to timestamp-based id.
    \item \texttt{--doc-id <string>} (optional): run pipeline for a single document id only.
    \item \texttt{--stages <list>} (optional): override selected stages (e.g., \texttt{baseline,cosine,preprocess,pass1,pass2}).
    \item \texttt{--strict} (optional): fail the run on schema validation errors; otherwise record failures and continue.
    \item \texttt{--log-level <level>} (optional): logging verbosity (e.g., INFO, WARNING, DEBUG).
\end{itemize}

\paragraph{Exit codes}
\begin{itemize}
    \item 0: success.
    \item 2: configuration error (invalid or missing config).
    \item 3: schema validation failure (strict mode).
    \item 4: unexpected runtime error.
\end{itemize}
